\SourceLeaf{208}\begin{center}
[\textsc{Bessel} \emph{to} \textsc{Gauss}. \emph{Königsberg in Prussia, 10 February 1829.}]
\end{center}

\{.\ .\ .\ .\ .\ .\ I would greatly regret it if you were deterred »by the clamor of the
Boeotians« from setting out your views on geometry.
What \textsc{Lambert} said and what \textsc{Schweikart} expressed orally
have made it clear to me that our geometry is incomplete and ought to receive a correction that is hypothetical and vanishes if the sum of the
angles of the plane triangle is $=180^{0}$. That would be the \textls{true}
geometry, and Euclidean geometry the \textls{practical} one, at least for figures on the
Earth. .\ .\ .\ .\ .\ .\}

\rule{4cm}{0.5pt}

\Needspace{16\baselineskip}
\begin{center}
\textsc{Gauss} to \textsc{Bessel}. Göttingen, 9 April 1830.
\end{center}

.\ .\ .\ .\ .\ .\ .\ I was truly glad at the ease with which you
engaged with my views on geometry, especially since so few
are open to them. I am deeply convinced that the theory of space has an entirely different place in our a priori knowledge from that of the pure
theory of magnitude. Our knowledge of the former is wholly lacking in the complete conviction of its necessity (and thus of its absolute
truth) that is characteristic of the latter. We must humbly admit
that, if number is \textls{merely} a product of our mind, space also has a reality outside
our mind, a reality whose laws we cannot fully prescribe a priori. .\ .\ .\ .\ .\ .\ .\ .

\rule{4cm}{0.5pt}

